\chapter{Quantization of the Electromagnetic Field}

\begin{remarkbox}
The scalar and Dirac fields taught us how particles and antiparticles arise from
mode expansions. The electromagnetic field adds a new feature: gauge redundancy.
A photon is not an excitation of four independent components of \(A_\mu\), but
of the physical transverse degrees of freedom of a gauge field.
\end{remarkbox}

\section{Why Gauge Fields Are Harder to Quantize}

The real scalar, complex scalar, and Dirac fields all had independent field
components that could be quantized after imposing the appropriate commutation or
anticommutation relations. The electromagnetic field is different. It is
described by a vector potential \(A_\mu\), but not all components of
\(A_\mu\) represent physical degrees of freedom.

The field strength is
\[
    F_{\mu\nu}=\partial_\mu A_\nu-\partial_\nu A_\mu.
\]
The free electromagnetic Lagrangian is
\[
    \mathcal{L}
    = -\frac{1}{4}F_{\mu\nu}F^{\mu\nu}.
\]
It is invariant under the gauge transformation
\[
    A_\mu\longrightarrow A_\mu+\partial_\mu\Lambda.
\]
The potentials \(A_\mu\) and \(A_\mu+\partial_\mu\Lambda\) describe the same
physical electromagnetic field. Therefore quantizing every component of
\(A_\mu\) as if it were an independent scalar field would overcount states.

\section{Physical and Redundant Degrees of Freedom}

A massless vector field in four spacetime dimensions has two physical
polarizations. This is already suggested classically by electromagnetic waves:
the electric and magnetic fields of a plane wave are transverse to the direction
of propagation.

The four components of \(A_\mu\) contain redundant information. Gauge freedom
removes one part, and the equations of motion impose a constraint. The result is
that only two transverse photon polarizations are physical.

This distinction is central:
\[
    \text{four components of } A_\mu
    \neq
    \text{four physical photon polarizations}.
\]
The task of quantization is to keep Lorentz covariance and locality under
control while ensuring that only the physical states contribute to observable
processes.

\section{Canonical Obstruction}

The canonical momentum conjugate to \(A_\mu\) is
\[
    \Pi^\mu
    = \frac{\partial\mathcal{L}}{\partial(\partial_0 A_\mu)}.
\]
For the time component one finds
\[
    \Pi^0=0.
\]
This means that \(A_0\) has no independent time derivative in the Lagrangian.
It is not an ordinary dynamical coordinate. Instead, it acts as a Lagrange
multiplier imposing Gauss's law.

This is the first signal that gauge fields require constraints or gauge fixing.
One cannot simply impose four independent canonical commutation relations for
all components of \(A_\mu\).

\section{Coulomb Gauge Quantization}

A physically transparent gauge choice is Coulomb gauge,
\[
    \nabla\cdot\vect{A}=0.
\]
For the free electromagnetic field, one may also set the nondynamical scalar
potential to zero after solving the constraint. The remaining field is the
transverse vector potential \(\vect{A}_T\).

The transverse field may be expanded as
\[
    \hat{\vect{A}}(t,\vect{x})
    = \sum_{\lambda=1}^{2}
      \int\frac{\dd^3k}{(2\pi)^3}
      \frac{1}{\sqrt{2|\vect{k}|}}
      \left(
          \vect{\epsilon}_{\lambda}(\vect{k})
          \hat a_{\lambda}(\vect{k})
          e^{-i|\vect{k}|t+i\vect{k}\cdot\vect{x}}
          +
          \vect{\epsilon}_{\lambda}^{\,*}(\vect{k})
          \hat a_{\lambda}^{\dagger}(\vect{k})
          e^{i|\vect{k}|t-i\vect{k}\cdot\vect{x}}
      \right).
\]
The polarization vectors are transverse:
\[
    \vect{k}\cdot\vect{\epsilon}_{\lambda}(\vect{k})=0.
\]
The index \(\lambda=1,2\) labels the two physical polarizations.

\section{Photon Creation and Annihilation Operators}

The photon ladder operators satisfy
\[
    [\hat a_{\lambda}(\vect{k}),
     \hat a_{\sigma}^{\dagger}(\vect{p})]
    = (2\pi)^3\delta_{\lambda\sigma}
      \delta^{(3)}(\vect{k}-\vect{p}),
\]
with all other commutators vanishing. The vacuum is defined by
\[
    \hat a_{\lambda}(\vect{k})|0\rangle=0
    \qquad \text{for all }\vect{k},\lambda.
\]
A one-photon state is
\[
    |\vect{k},\lambda\rangle
    = \hat a_{\lambda}^{\dagger}(\vect{k})|0\rangle.
\]
Its energy is
\[
    \omega_{\vect{k}}=|\vect{k}|,
\]
which is the massless dispersion relation.

The normal-ordered Hamiltonian is
\[
    :\hat H:
    = \sum_{\lambda=1}^{2}
      \int\frac{\dd^3k}{(2\pi)^3}\,
      |\vect{k}|\,
      \hat a_{\lambda}^{\dagger}(\vect{k})
      \hat a_{\lambda}(\vect{k}).
\]
The momentum operator is
\[
    :\hat{\vect{P}}:
    = \sum_{\lambda=1}^{2}
      \int\frac{\dd^3k}{(2\pi)^3}\,
      \vect{k}\,
      \hat a_{\lambda}^{\dagger}(\vect{k})
      \hat a_{\lambda}(\vect{k}).
\]

\section{Polarization Vectors}

For each nonzero momentum \(\vect{k}\), one chooses two transverse polarization
vectors. They obey
\[
    \vect{k}\cdot\vect{\epsilon}_{\lambda}(\vect{k})=0,
    \qquad
    \vect{\epsilon}_{\lambda}^{\,*}(\vect{k})
    \cdot
    \vect{\epsilon}_{\sigma}(\vect{k})
    = \delta_{\lambda\sigma}.
\]
Linear polarizations are useful for many classical comparisons. Circular
polarizations are useful for describing helicity. They are built from two linear
polarizations by
\[
    \vect{\epsilon}_{\pm}
    = \frac{1}{\sqrt{2}}
      \left(\vect{\epsilon}_1\pm i\vect{\epsilon}_2\right).
\]
The labels \(+\) and \(-\) correspond to the two helicity states of the photon.

\section{Photon States and Helicity}

A massive spin-one particle would have three spin polarizations. A photon is
massless and gauge invariant, so only two physical helicities remain. A
one-photon state may be labelled as
\[
    |\vect{k},+\rangle,
    \qquad
    |\vect{k},-\rangle.
\]
These are helicity eigenstates. Helicity is the projection of angular momentum
onto the direction of motion. For a massless particle, helicity is Lorentz
invariant under proper orthochronous Lorentz transformations.

The absence of a longitudinal physical photon is not an accident. The
longitudinal component of \(\vect{A}\) is gauge redundancy, not a separate
particle.

\section{Covariant Quantization}

Coulomb gauge makes the physical degrees of freedom visible, but it hides
manifest Lorentz covariance. For perturbative QFT it is often more convenient to
use a covariant gauge-fixing term,
\[
    \mathcal{L}_{\mathrm{gf}}
    = -\frac{1}{2\xi}(\partial_\mu A^\mu)^2.
\]
The gauge-fixed Lagrangian is
\[
    \mathcal{L}
    = -\frac{1}{4}F_{\mu\nu}F^{\mu\nu}
      -\frac{1}{2\xi}(\partial_\mu A^\mu)^2.
\]
The parameter \(\xi\) labels a family of gauges. The choice \(\xi=1\) is called
Feynman gauge.

Covariant quantization introduces unphysical components into intermediate
calculations. These components must not appear in physical observables. In QED,
current conservation ensures that unphysical polarizations cancel from
scattering amplitudes. In non-Abelian gauge theory, the corresponding statement
requires ghosts and BRST symmetry, which will be developed later.

\section{Gauge Fixing and the Photon Propagator}

In covariant gauges, the momentum-space photon propagator is
\[
    D_{\mu\nu}(k)
    = \frac{-i}{k^2+i\epsilon}
      \left(
          \eta_{\mu\nu}
          - (1-\xi)\frac{k_\mu k_\nu}{k^2+i\epsilon}
      \right).
\]
In Feynman gauge, this simplifies to
\[
    D_{\mu\nu}(k)
    = \frac{-i\eta_{\mu\nu}}{k^2+i\epsilon}.
\]
The pole at \(k^2=0\) reflects the masslessness of the photon. The numerator
carries the vector structure of the field.

The gauge-dependent part proportional to \(k_\mu k_\nu\) does not change
physical QED amplitudes coupled to conserved currents. This is the practical
reason covariant gauges are useful: they simplify calculations without changing
observable predictions.

\section{Unphysical Polarizations}

Covariant quantization treats \(A_\mu\) as a four-component field during
intermediate steps. This introduces time-like and longitudinal polarizations
that are not physical photon states. The physical Hilbert space is obtained only
after imposing the appropriate condition that removes the unphysical sector.

In elementary QED calculations, this removal is often hidden inside Ward
identities and current conservation. If an external photon polarization vector is
shifted by a multiple of its momentum,
\[
    \epsilon_\mu(k)\longrightarrow \epsilon_\mu(k)+c\,k_\mu,
\]
physical amplitudes are unchanged. This is the amplitude-level expression of
gauge redundancy.

\section{Gauge-Invariant Observables}

The potential \(A_\mu\) is not gauge invariant. The field strength
\(F_{\mu\nu}\) is gauge invariant, and so are the electric and magnetic fields.
Observables must ultimately be built from gauge-invariant quantities or from
gauge-invariant combinations of charged fields and gauge fields.

Examples include
\[
    F_{\mu\nu}(x),
    \qquad
    F_{\mu\nu}(x)F^{\mu\nu}(x),
\]
and, in more advanced settings, Wilson line and Wilson loop operators. The
important lesson at this stage is that the photon is described by a gauge
potential, but physical information cannot depend on the redundant gauge choice.

\section{What Is a Photon?}

In free QED, a photon is a one-quantum excitation of a transverse mode of the
electromagnetic field:
\[
    |\vect{k},\lambda\rangle
    = \hat a_{\lambda}^{\dagger}(\vect{k})|0\rangle,
    \qquad
    \lambda=1,2.
\]
Equivalently, one may label the two states by helicity \(\lambda=\pm\). The
photon is massless, has energy \(|\vect{k}|\), and carries no electric charge.

This definition is sharp in the free theory and in scattering theory with
asymptotic photon states. In interacting QED, soft photons and infrared effects
make the exact particle interpretation more subtle. Those subtleties will be
revisited when discussing scattering and loop corrections.

\begin{summarybox}
The electromagnetic field is the first gauge field in the course. Gauge
redundancy means that not all components of \(A_\mu\) are physical. In Coulomb
gauge the two transverse photon polarizations are explicit. In covariant gauges
one keeps Lorentz covariance at the cost of introducing unphysical components,
which cancel from physical observables. A photon is a transverse, massless
quantum of the electromagnetic field.
\end{summarybox}

\section{Chapter Checklist}

After reading this chapter, one should be able to:
\begin{itemize}
    \item explain why gauge fields are harder to quantize than scalar fields;
    \item state the gauge redundancy of \(A_\mu\);
    \item explain why the photon has two physical polarizations;
    \item write the Coulomb-gauge mode expansion of the transverse vector
    potential;
    \item construct one-photon states;
    \item distinguish linear polarization from helicity;
    \item write the covariant gauge-fixing term;
    \item write the photon propagator in a covariant gauge;
    \item explain why unphysical polarizations do not appear as external photon
    states;
    \item identify gauge-invariant electromagnetic observables.
\end{itemize}
